\section{本章小结}
\label{sec:rvm-summary}

本章从一张概率表出发，问怎样用少数几个数概括它。期望是概率质量的重心，按概率加权平均取值；方差是绕重心的转动惯量，按概率加权平均偏差的平方；更高阶的矩加重远处偏差的权重，标准化后的三阶、四阶矩描述偏斜方向与$Z^{2}$的波动。奖金表与密度$2x$贯穿全章，表~\ref{tab:rvm-summary}汇总了它们的各项数值。

\begin{BookInlineTable}
 \centering
 \caption{两个贯穿例子的数值特征。}
 \label{tab:rvm-summary}
 \begin{tabular}{lcc}
  \toprule
  数值特征 & 奖金表 & 密度$2x$\\
  \midrule
  均值$\mu$ & $2$ & $2/3$\\
  方差$\sigma^{2}$ & $5$ & $1/18$\\
  标准差$\sigma$ & $2.24$ & $0.236$\\
  三阶中心矩$c_3$ & $7$ & $-1/135$\\
  四阶中心矩$c_4$ & $51$ & $1/135$\\
  偏度$\gamma_1$ & $0.626$ & $-0.566$\\
  峰度$\kappa$ & $2.04$ & $2.4$\\
  \bottomrule
 \end{tabular}
\end{BookInlineTable}

每个数回答一个特定的问题，也需要特定的条件，表~\ref{tab:rvm-summary-roles}把它们与第~\ref{chap:random-variables}章的分位数放在一起比较。均值与中位数都描述位置，前者使平方误差最小、受远处取值牵动，后者使绝对误差最小、只看左右质量的分配。方差与标准差描述波动尺度，标准差与变量同单位。偏度与峰度不随单位改变，却只是两个数，零偏度不保证对称，峰度也不能确定峰形或尾部规律。

\begin{BookInlineTable}
 \centering
 \caption{各数值特征回答的问题、存在条件与在$aX+b$（$a\neq0$）下的变化。偏度与峰度还要求$\sigma>0$。}
 \label{tab:rvm-summary-roles}
 \begin{tabular}{llll}
  \toprule
  特征 & 回答的问题 & 存在条件 & 在$aX+b$下\\
  \midrule
  均值 & 重心在哪里 & $\mathbb{E}\lvert X\rvert<\infty$ & $a\mu+b$\\
  中位数 & 左右各半的位置 & 总存在 & 随之变换\\
  方差 & 绕重心散开多少 & $\mathbb{E}[X^{2}]<\infty$ & $a^{2}\sigma^{2}$\\
  标准差 & 同单位的波动 & $\mathbb{E}[X^{2}]<\infty$ & $\lvert a\rvert\sigma$\\
  $k$阶中心矩 & 偏差的$k$次平均 & $\mathbb{E}\lvert X\rvert^{k}<\infty$ & $a^{k}c_k$\\
  偏度 & 偏斜方向 & $\mathbb{E}\lvert X\rvert^{3}<\infty$ & $a<0$时变号\\
  峰度 & $Z^{2}$的波动 & $\mathbb{E}[X^{4}]<\infty$ & 不变\\
  \bottomrule
 \end{tabular}
\end{BookInlineTable}

存在条件是使用这些数的前提。非负变量的期望总有意义但可能为$+\infty$；有符号变量要分别检查正部与负部，两者都无穷时期望不存在，对称也不能补救。高阶绝对矩有限蕴含低阶绝对矩有限，Lyapunov不等式给出了定量的比较；反方向不成立，幂律尾部的密度$\alpha x^{-\alpha-1}$只有低于$\alpha$阶的矩。一切展开公式和对称性论证都要先确认所涉最高阶的绝对矩有限。

运算规则同样有边界。期望的线性性对可积变量的有限线性组合成立，不需要独立性，信封问题说明它能绕开复杂的依赖结构。乘积的期望一般不能拆开，无穷求和与取期望的交换需要额外条件。函数的期望可以直接用原变量的分布计算，不必求出新变量的分布，但要对每个函数检查绝对可积。凸函数满足Jensen不等式，变量越分散，“先平均再变换”与“先变换再平均”相差越大。方差的展开式便于手算，数值计算中则应先减去均值。

数值特征是分布的概括，不是替代。均值和方差相同的概率表形状可以完全不同；对任意给定的$k$，都有前$k$阶矩相同的不同密度。尾概率的估计、和的方差以及用变换识别分布，分别在第~\ref{chap:tail-probabilities}章、第~\ref{chap:random-vectors}章与第~\ref{chap:weak-convergence-of-distributions}章展开。下一章把本章的工具用于常见分布，在每个模型的试验机制和参数下计算均值与方差，并检查它们是否存在。
